% Figura: flujo de un campo a través de una superficie. (a) superficie
% ondulada inmersa en un campo; (b) elemento de área dirigido y componente normal.
\begin{figure}[ht]
\centering
\begin{minipage}[t]{0.56\textwidth}\centering
\begin{tikzpicture}[baseline=(current bounding box.north)]
  \begin{axis}[hide axis, width=1.12\textwidth, height=0.95\textwidth, view={-28}{40},
      xmin=-3.2, xmax=3.2, ymin=-2.4, ymax=2.4, zmin=-1.6, zmax=2.2,
      colormap name=verdeamarillo, point meta min=-0.9, point meta max=0.9,
      /pgf/number format/use comma, clip=false]
    % campo por debajo de la superficie
    \addplot3[quiver={u=0.34, v=0.18, w=0.80}, -{Stealth[length=3pt]}, line width=0.7pt,
              acento, domain=-3.0:3.0, domain y=-2.2:2.2, samples=6, samples y=5]
      {-1.2};
    % superficie
    \addplot3[surf, shader=interp, samples=28, samples y=22, domain=-2.5:2.5, domain y=-1.8:1.8,
              opacity=0.93, mesh/ordering=x varies]
      {0.85*sin(deg(1.25*x))*cos(deg(0.5*y))+0.12*y};
    % campo por encima de la superficie
    \addplot3[quiver={u=0.34, v=0.18, w=0.80}, -{Stealth[length=3pt]}, line width=0.7pt,
              acento, domain=-3.0:3.0, domain y=-2.2:2.2, samples=6, samples y=5]
      {1.0};
    \node[font=\small, text=acento] at (axis cs:2.6,0.3,1.9) {$\vc{A}$};
    \node[font=\small, text=black!80, fill=white, fill opacity=0.7, text opacity=1, inner sep=1pt] at (axis cs:-2.0,-1.2,-0.5) {$S$};
  \end{axis}
\end{tikzpicture}\\[2pt]{\small\scshape (a)}
\end{minipage}\hfill
\begin{minipage}[t]{0.40\textwidth}\centering
\begin{tikzpicture}[>=Stealth, x=1cm, y=1cm, baseline=(current bounding box.north)]
  \shade[top color=azulclaro!55, bottom color=azulfondo, shading angle=0]
    (0,0) -- (3.4,0) -- (2.7,1.1) -- (-0.7,1.1) -- cycle;
  \draw[grispizarra!70, line width=0.4pt] (0,0) -- (3.4,0) -- (2.7,1.1) -- (-0.7,1.1) -- cycle;
  \coordinate (O) at (1.35,0.55);
  \draw[->, line width=1.6pt, azulacero] (O) -- ++(0,2.1) node[left, text=azulacero] {$\dS$};
  \draw[->, line width=1.6pt, rojoteo] (O) -- ++(2.0,1.55) node[right, text=rojoteo] {$\vc{A}$};
  % componente normal
  \draw[densely dashed, grispizarra] (3.35,2.1) -- (1.35,2.1);
  \draw[line width=3pt, azulclaro, opacity=0.9] (O) -- ++(0,1.55);
  \node[font=\footnotesize, text=azulprof, anchor=east] at (1.25,1.5) {$\vc{A}\cdot\un$};
  \draw[grispizarra, line width=0.5pt] (1.35,1.15) arc (90:37.8:0.6);
  \node[font=\small] at (1.85,1.25) {$\theta$};
  \fill (O) circle (1.4pt);
\end{tikzpicture}\\[2pt]{\small\scshape (b)}
\end{minipage}
\caption{Flujo de un campo vectorial $\vc{A}$ a través de una superficie $S$.
(a)~Un campo (flechas rojas) atraviesa una superficie curva; el flujo
contabiliza solo la parte del campo que corta a la superficie. (b)~En un
entorno de cada punto la superficie es plana: $\dS=\un\,\dsc$ es un vector de
módulo igual al área del parche y perpendicular a él, y el flujo
elemental es $\vc{A}\cdot\dS=(\vc{A}\cdot\un)\,\dsc=\norm{\vc{A}}\cos\theta\,\dsc$.}
\label{fig:flujo}
\end{figure}
